\documentclass[12pt]{article}
\usepackage[utf8]{inputenc}
\usepackage{amsmath, amssymb} % Math packages
\usepackage{geometry}
\usepackage{hyperref} % Clickable links
\usepackage{graphicx} % Images

\geometry{a4paper, margin=1in}

\title{Sistemi operativi appunti}
\author{Matteo}
\date{\today}

\begin{document}

\maketitle

\section{secondo semestre}

sistema operativo: cio che crea l'astrazione dei processi

librerie sono gruppi di funzioni che servono a rendere piu facile l'uso del sistema, non sono necessarie.

per il sistema operativo, le system call sono cio che forniscono astrazioni, come escuzzione di programmi o accesso al file system.

\subsection{storia dei sistemi operativi}

4 generazioni + la gen 0
\begin{enumerate}
\item gen 0= macchine meccaniche
\item gen 1= macchine a valvole programmate con tavole di commutazione. guasti molto frequenti.
\item gen 2= vengono introdotti i transistor, macchine piu efficaci ed economiche. programmati con i primi linguaggi ad alto livello.
\item gen 3= arrivano i circuiti integrati, cioe tanti transistor su un unico componente, i transistor sono poi collegati. Nasce multiprogrammazione. nasce la necessita di gestone della memoria
\item gen 4= un singolo elaboratore ora possiede piu unita di elaborazione, nascono i personal computer, nascono problemi di sicurezza.

\end{enumerate}

sistemi real-time= sistemi in cui la correttezza dell'elaborazione non viene solo misurata dalla correttezza della risposta ma anche dal tempo impiegato(devono rispondere entro un dato periodo)




\end{document}
